\documentclass[11pt,a4paper]{article}
\usepackage[T1]{fontenc}\usepackage[utf8]{inputenc}\usepackage[czech]{babel}
\usepackage{lmodern,microtype,amsmath,amssymb,geometry,enumitem}
\geometry{top=18mm,bottom=18mm,left=23mm,right=23mm}
\pagestyle{empty}\setlength{\parindent}{0pt}
\begin{document}
\begin{center}{\Large\bfseries Matematická únikovka -- učitelský klíč}\end{center}
\begin{enumerate}[label=\arabic*.]
\item Sudá $A=\{2;4;6;8\}$, dělitelná třemi $B=\{3;6;9\}$. Právě jednu z vlastností splňuje $A\triangle B=\{2;3;4;8;9\}$, tedy \textbf{5}.
\item Pro $x\ne3$ platí $\frac{x^2-9}{x-3}=x+3=7$, takže $x=\textbf{4}$.
\item Směrnice $a=(4-8)/(3-1)=-2$, předpis $f(x)=-2x+10$, tedy $f(2)=\textbf{6}$.
\item Sečtením rovnic získáme $3x=12$, tedy $x=4$ a $y=\textbf{5}$.
\item Z $|x-4|=3$ plyne $x=1$ nebo $x=7$, součet je \textbf{8}.
\item $x^2-5x+6=(x-2)(x-3)$, kořeny $2,3$ a jejich součin je \textbf{6}.
\end{enumerate}
\bigskip\begin{center}{\Huge\bfseries 546586}\end{center}
\medskip
\textbf{Metodická poznámka:} V posledním stanovišti si žáci mohou využít i Viètovy vztahy. Kontroluj, zda při krácení v úloze 2 uvedli $x\ne3$.
\end{document}
